\begin{questionS}\label{Q:consumerM}
Consider the following producer-consumer problem.  Several producers each
produce one piece of integer data at a time, which they place into a buffer
using a method \SCALA{put}.  These are accumulated into an array of size~|n|.
A single consumer can receive that array, once it is full, using the method
\SCALA{get}.  Thus the desired functionality is captured by the following
trait.
%
\begin{scala}
trait ProducerConsumer{
  /** Add £x£ to the buffer, waiting until it is not full. */
  def put(x: Int): Unit

  /** Get the contents of the buffer, once it is full. */
  def get: Array[Int]
}
\end{scala}
%
Note that there is a single consumer that calls the |get| operation. 

Implement the buffer as a monitor.  The consumer should be blocked until the
array is full, and the producers should be blocked once the array is full.
Give your implementation the following signature.
%
\begin{scala}
class ProducerConsumerMonitor(n: Int) extends ProducerConsumer{
  require(n >= 1)
  ...
}
\end{scala}

Write testing code for your implementation, using the linearizability testing
framework.  Remember that the results of the concurrent and sequential
operations need to be comparable using ``|==|'', which won't be the case for
|Array[Int]|.  I suggest you cast the result of the concurrent |get| operation
to |List[Int]| using the |_.toList| function, and arrange for the
corresponding sequential operation to return a result of the same type.
Remember also that the sequential specification object needs to be immutable,
and ideally comparable using the ``|==|'' operation.
\end{questionS}

%%%%%

\begin{answerS}
My code is below.
%
\begin{scala}
class MonitorProducerConsumer(n: Int) extends ProducerConsumer{
  require(n >= 1)

  /** Array to hold the data. */
  private var a = new Array[Int](n)

  /** Number of items added to £a£ so far. */
  private var count = 0

  private val lock = new Lock

  /** Condition for signalling that the old buffer has been taken. */
  private val notFull = lock.newCondition

  /** Condition for signalling that the buffer has been filled. */
  private val full = lock.newCondition

  /** Add £x£ to the buffer, waiting until it is not full. */
  def put(x: Int) = lock.mutex{
    notFull.await(count < n) // Wait for space (1).
    a(count) = x; count += 1
    if(count == n) full.signal() // Signal to a consumer at (2).
  }

  /** Get the contents of the buffer, once it is full. */
  def get: Array[Int] = lock.mutex{
    if(count < n) full.await() // Wait for the buffer to be filled (2).
    val result = a; a = new Array[Int](n); count = 0;
    notFull.signalAll() // Signal to all the producers at (1).
    result    
  }
\end{scala}
%
Note that the producer has to recheck whether the buffer full each time it is
awoken, since there are multiple producers, and it is possible that other
producers have filled the new buffer in the meantime.  However, there is no
need for the consumer to recheck each time it is awoken, since there is only
one consumer.

For testing, we take the sequential specification object to be a |List[Int]|.
The last part of the question hints why using an |Array[Int]| would be less
good (it would require cloning the array, and objects wouldn't be comparable
using~``|==|'').  Following the question, we  arrange for the sequential get
operation to return a |List[Int]|, which is easy given the choice of
sequential specification type. 
%
\begin{scala}
  /** Sequential specification type.  */
  type S = List[Int]

  def seqPut(x: Int)(buffer: S): (Unit, S) = {
    require(buffer.length < n); ((), buffer :+ x)
  }

  def seqGet(buffer: S): (List[Int], S) = {
    require(buffer.length == n); (buffer, List[Int]())
  }
\end{scala}
%
Each consumer needs to cast the result of the |get| to a |List|.
%
\begin{scala}
  /** A worker thread. Thread 0 is the consumer. */
  def worker(me: Int, log: LinearizabilityLog[S,ProducerConsumer]) = {
    val random = new scala.util.Random
    if(me == 0) for(i <- 0 until cIters) log(_.get.toList, "get", seqGet)
    else 
      for(i <- 0 until pIters){
        val x = random.nextInt(maxValue); log(_.put(x), s"put($x)", seqPut(x))
      }
  }
\end{scala}
%
I arranged for |numProducers| producers, each performing |pIters| iterations
per run, and a single consumer performing |cIters| iterations per run, where
$\sm{pIters} \times \sm{numProducers} = \sm n \times \sm{cIters}$.
The rest of the testing is standard.
\end{answerS}
